\documentclass[oneside]{book}
\usepackage{xeCJK}
\usepackage{atbegshi}
\usepackage{jiazhu}

\defaultCJKfontfeatures{RawFeature={vertical:+vert:+vhal}}
\setCJKmainfont{Noto Serif CJK SC}
\newCJKfontfamily{\jiazhufont}{Noto Serif CJK SC}[FakeStretch=1.6]

\jiazhuset{
  format=\jiazhufont,
  ratio=0.625,
  beforeskip=0pt,
  afterskip=0pt,
  shortcut=【】
}
\parindent=0pt
\AtBeginShipout{
  \special { pdf:put @pages<</Rotate 90>>}
}

\setlength{\textwidth}{400pt}

\begin{document}
\fontsize{20pt}{30pt}\selectfont

學之書古之大學所以教人之法也蓋自天降生民則既莫不與之以仁義禮
智之性矣然其氣質之【稟或不能齊是以不能皆有以知其性之所】有而全之也
一有聰明睿智能盡其性者出於其閑則天必命之以爲億兆之君師使之治而
教之以複其性此伏羲神農黃帝堯舜所以繼天立極而司徒之職典樂之官所
由設也

三代之隆其【法寖備然後王宮國都以及閭巷莫不有學人生八歲則自王公以
下至於庶人之子弟皆入小學而教之以灑掃應對進退之節禮樂射禦書數之
文及其十有五年則自天子之元子眾子以至公卿大夫元士之適子與凡民之
俊秀皆入大學而教之以窮理正心修己治人之道此又學校之教大小之節所
以分也】

夫以學校之設其廣如此教之之術其次第節目之詳又如此而其所以爲教則
又皆本之人君躬行心得之餘不待求之民生日用彝倫之外是以當世之人無
不學其學焉者無不有以知其性分之所固有職分之所當爲而各俛焉以盡其
力此古昔盛時所以治隆於上俗美於下而非後世之所能及也
\end{document}
